\textbf{Phần 1. Trắc nghiệm lựa chọn}
\setcounterpageref{ex}{0}

\begin{ex} %[5H5V26]
Một hình thang có diện tích bằng $105\text{ cm}^2$ và chiều cao là $7\text{ cm}$. Trung bình cộng độ dài hai đáy của hình thang đó là
\choice
{$30\text{ cm}$}
{\True $15\text{ cm}$}
{$10\text{ cm}$}
{$20\text{ cm}$}
\loigiai{
Trung bình cộng hai đáy bằng diện tích chia cho chiều cao:
$$105 : 7 = 15\text{ (cm)}.$$
\textbf{Đáp án của câu hỏi là B}
}
\end{ex}

\begin{ex} %[5H5V26]
Một mảnh ruộng hình thang có đáy lớn $28\text{ m}$, đáy bé $18\text{ m}$. Chiều cao của mảnh ruộng bằng trung bình cộng độ dài hai đáy. Diện tích của mảnh ruộng đó là
\choice
{$460\text{ m}^2$}
{\True $529\text{ m}^2$}
{$512\text{ m}^2$}
{$540\text{ m}^2$}
\loigiai{
Chiều cao mảnh ruộng là:
$(28 + 18) : 2 = 23\text{ (m)}$.\\
Diện tích mảnh ruộng là:
$23 \times 23 = 529\text{ (m}^2\text{)}$.\\
\textbf{Đáp án của câu hỏi là B}
}
\end{ex}

\begin{ex} %[5H5V26]
Một hình thang có diện tích $180\text{ m}^2$, chiều cao $12\text{ m}$. Biết đáy lớn hơn đáy bé $6\text{ m}$. Độ dài đáy lớn của hình thang đó là
\choice
{$12\text{ m}$}
{\True $18\text{ m}$}
{$24\text{ m}$}
{$15\text{ m}$}
\loigiai{
Tổng độ dài hai đáy là:
$180 \times 2 : 12 = 30\text{ (m)}$.\\
Độ dài đáy lớn là:
$(30 + 6) : 2 = 18\text{ (m)}$.\\
\textbf{Đáp án của câu hỏi là B}
}
\end{ex}

\begin{ex} %[5H5V26]
Một thửa ruộng hình thang có diện tích $360\text{ m}^2$. Nếu mở rộng đáy lớn thêm $5\text{ m}$ thì diện tích thửa ruộng tăng thêm $40\text{ m}^2$. Tổng độ dài hai đáy ban đầu của thửa ruộng là
\choice
{$40\text{ m}$}
{\True $45\text{ m}$}
{$50\text{ m}$}
{$36\text{ m}$}
\loigiai{
Phần diện tích tăng thêm là hình tam giác có đáy $5\text{ m}$ và chiều cao bằng chiều cao hình thang.\\
Chiều cao thửa ruộng là:
$40 \times 2 : 5 = 16\text{ (m)}$.\\
Tổng độ dài hai đáy ban đầu là:
$360 \times 2 : 16 = 45\text{ (m)}$.\\
\textbf{Đáp án của câu hỏi là B}
}
\end{ex}

\begin{ex} %[5H5V26]
Một hình thang vuông có chiều cao bằng đáy bé và bằng $12\text{ cm}$. Đáy lớn gấp đôi đáy bé. Diện tích của hình thang vuông đó là
\choice
{$144\text{ cm}^2$}
{\True $216\text{ cm}^2$}
{$288\text{ cm}^2$}
{$432\text{ cm}^2$}
\loigiai{
Đáy lớn của hình thang vuông là:
$12 \times 2 = 24\text{ (cm)}$.\\
Diện tích của hình thang vuông là:
$(12 + 24) \times 12 : 2 = 216\text{ (cm}^2\text{)}$.\\
\textbf{Đáp án của câu hỏi là B}
}
\end{ex}

\begin{ex} %[5H5V26]
Một mảnh đất hình thang có đáy lớn $35\text{ m}$, đáy bé bằng $\dfrac{3}{5}$ đáy lớn, chiều cao kém đáy bé $5\text{ m}$. Diện tích mảnh đất đó là
\choice
{$420\text{ m}^2$}
{\True $448\text{ m}^2$}
{$480\text{ m}^2$}
{$896\text{ m}^2$}
\loigiai{
Đáy bé mảnh đất là:
$35 \times \dfrac{3}{5} = 21\text{ (m)}$.\\
Chiều cao mảnh đất là:
$21 - 5 = 16\text{ (m)}$.\\
Diện tích mảnh đất là:
$(35 + 21) \times 16 : 2 = 448\text{ (m}^2\text{)}$.\\
\textbf{Đáp án của câu hỏi là B}
}
\end{ex}

\begin{ex} %[5H5V26]
Một thửa ruộng hình thang có đáy lớn $26\text{ m}$, đáy bé $14\text{ m}$ và chiều cao $15\text{ m}$. Cứ mỗi $100\text{ m}^2$ thu hoạch được $70\text{ kg}$ thóc. Khối lượng thóc thu hoạch được trên thửa ruộng đó là
\choice
{$180\text{ kg}$}
{\True $210\text{ kg}$}
{$240\text{ kg}$}
{$280\text{ kg}$}
\loigiai{
Diện tích thửa ruộng là:
$(26 + 14) \times 15 : 2 = 300\text{ (m}^2\text{)}$.\\
Khối lượng thóc thu hoạch được là:
$300 : 100 \times 70 = 210\text{ (kg)}$.\\
\textbf{Đáp án của câu hỏi là B}
}
\end{ex}

\begin{ex} %[5H5V26]
Nếu tăng cả hai đáy của một hình thang lên gấp 2 lần và giữ nguyên chiều cao thì diện tích hình thang sẽ
\choice
{Tăng gấp 4 lần}
{\True Tăng gấp 2 lần}
{Không thay đổi}
{Tăng gấp 8 lần}
\loigiai{
Diện tích hình thang ban đầu: $S = \dfrac{(a + b) \times h}{2}$.\\
Khi tăng hai đáy lên 2 lần: $S' = \dfrac{(2a + 2b) \times h}{2} = 2 \times S$.\\
Vậy diện tích tăng gấp 2 lần.\\
\textbf{Đáp án của câu hỏi là B}
}
\end{ex}

\begin{ex} %[5H5V26]
Một mảnh vườn hình thang có đáy bé $16\text{ m}$, đáy lớn $24\text{ m}$. Người ta thu hẹp đáy lớn bớt $4\text{ m}$ để làm đường đi thì diện tích mảnh vườn giảm đi $28\text{ m}^2$. Chiều cao của mảnh vườn là
\choice
{$12\text{ m}$}
{\True $14\text{ m}$}
{$16\text{ m}$}
{$18\text{ m}$}
\loigiai{
Phần diện tích giảm đi là tam giác có đáy $4\text{ m}$ và chiều cao bằng chiều cao hình thang.\\
Chiều cao mảnh vườn là:
$28 \times 2 : 4 = 14\text{ (m)}$.\\
\textbf{Đáp án của câu hỏi là B}
}
\end{ex}

\begin{ex} %[5H5V26]
Một hình thang có đáy lớn $40\text{ cm}$, đáy bé bằng $60\%$ đáy lớn. Chiều cao bằng một nửa tổng hai đáy. Diện tích của hình thang là
\choice
{$512\text{ cm}^2$}
{\True $1\,024\text{ cm}^2$}
{$2\,048\text{ cm}^2$}
{$960\text{ cm}^2$}
\loigiai{
Đáy bé của hình thang là:
$40 \times 60\% = 24\text{ (cm)}$.\\
Tổng hai đáy là:
$40 + 24 = 64\text{ (cm)}$.\\
Chiều cao hình thang là:
$64 : 2 = 32\text{ (cm)}$.\\
Diện tích hình thang là:
$64 \times 32 : 2 = 1\,024\text{ (cm}^2\text{)}$.\\
\textbf{Đáp án của câu hỏi là B}
}
\end{ex}

\begin{ex} %[5H5V26]
Một mặt bàn hình thang có đáy lớn $1{,}4\text{ m}$, đáy bé $0{,}8\text{ m}$ và chiều cao $0{,}6\text{ m}$. Chi phí làm mặt bàn là $500\,000$ đồng cho mỗi mét vuông. Tổng chi phí làm mặt bàn là
\choice
{$300\,000\text{ đồng}$}
{\True $330\,000\text{ đồng}$}
{$360\,000\text{ đồng}$}
{$660\,000\text{ đồng}$}
\loigiai{
Diện tích mặt bàn là:
$(1{,}4 + 0{,}8) \times 0{,}6 : 2 = 0{,}66\text{ (m}^2\text{)}$.\\
Chi phí làm mặt bàn là:
$0{,}66 \times 500\,000 = 330\,000\text{ (đồng)}$.\\
\textbf{Đáp án của câu hỏi là B}
}
\end{ex}

\begin{ex} %[5H5V26]
Cho hình thang $ABCD$ (đáy bé $AB$, đáy lớn $CD$). Hai đường chéo $AC$ và $BD$ cắt nhau tại điểm $O$. Khẳng định nào sau đây luôn đúng?
\choice
{Diện tích tam giác $AOB$ bằng diện tích tam giác $COD$}
{\True Diện tích tam giác $AOD$ bằng diện tích tam giác $BOC$}
{Diện tích tam giác $ABC$ lớn hơn diện tích tam giác $ABD$}
{Diện tích tam giác $ACD$ bé hơn diện tích tam giác $BCD$}
\loigiai{
Hai tam giác $ABD$ và $ABC$ có chung đáy $AB$ và cùng chiều cao bằng chiều cao hình thang nên $S_{ABD} = S_{ABC}$.\\
Bớt đi phần chung $S_{AOB}$, ta được:
$$S_{AOD} = S_{BOC}.$$
\textbf{Đáp án của câu hỏi là B}
}
\end{ex}

\textbf{Phần 2. Trắc nghiệm đúng sai}
\setcounterpageref{ex}{0}

\begin{ex} %[5H5V26]
Cho hình thang có độ dài đáy lớn là $32\text{ cm}$, đáy bé là $18\text{ cm}$ và chiều cao $14\text{ cm}$.
\choiceTF
{\True Tổng độ dài hai đáy của hình thang bằng $50\text{ cm}$}
{\True Trung bình cộng độ dài hai đáy của hình thang bằng $25\text{ cm}$}
{Diện tích của hình thang đó bằng $700\text{ cm}^2$}
{\True Nếu giữ nguyên hai đáy và tăng chiều cao thêm $2\text{ cm}$ thì diện tích hình thang tăng thêm $50\text{ cm}^2$}
\loigiai{
\begin{itemchoice}
\item Đúng. Tổng hai đáy: $32 + 18 = 50\text{ (cm)}$.
\item Đúng. Trung bình cộng: $50 : 2 = 25\text{ (cm)}$.
\item Sai. Diện tích hình thang: $25 \times 14 = 350\text{ (cm}^2\text{)} 
eq 700\text{ cm}^2$.
\item Đúng. Diện tích tăng thêm là: $50 \times 2 : 2 = 50\text{ (cm}^2\text{)}$.
\end{itemchoice}
}
\end{ex}

\begin{ex} %[5H5V26]
Một thửa ruộng hình thang có đáy lớn $45\text{ m}$, đáy bé $35\text{ m}$ và chiều cao $24\text{ m}$.
\choiceTF
{\True Diện tích của thửa ruộng là $960\text{ m}^2$}
{\True Người ta dành $25\%$ diện tích thửa ruộng để đào ao thả cá, diện tích ao là $240\text{ m}^2$}
{Phần diện tích còn lại dùng để trồng lúa là $700\text{ m}^2$}
{\True Diện tích trồng lúa gấp 3 lần diện tích đào ao}
\loigiai{
\begin{itemchoice}
\item Đúng. Diện tích thửa ruộng: $(45 + 35) \times 24 : 2 = 960\text{ (m}^2\text{)}$.
\item Đúng. Diện tích đào ao: $960 \times 25\% = 240\text{ (m}^2\text{)}$.
\item Sai. Diện tích trồng lúa là: $960 - 240 = 720\text{ (m}^2\text{)} 
eq 700\text{ m}^2$.
\item Đúng. Tỉ số diện tích trồng lúa và đào ao là: $720 : 240 = 3$ (lần).
\end{itemchoice}
}
\end{ex}

\textbf{Phần 3. Trắc nghiệm trả lời ngắn}
\setcounterpageref{ex}{0}

\begin{ex} %[5H5V26]
Một hình thang có diện tích $156\text{ cm}^2$ và chiều cao $12\text{ cm}$. Tổng độ dài hai đáy của hình thang đó là bao nhiêu xăng-ti-mét?
\shortans{26}
\loigiai{
Tổng độ dài hai đáy của hình thang là:
$156 \times 2 : 12 = 26\text{ (cm)}$.
}
\end{ex}

\begin{ex} %[5H5V26]
Một hình thang có diện tích $120\text{ m}^2$, độ dài đáy lớn là $14\text{ m}$ và đáy bé là $10\text{ m}$. Chiều cao của hình thang đó bằng bao nhiêu mét?
\shortans{10}
\loigiai{
Chiều cao của hình thang là:
$120 \times 2 : (14 + 10) = 240 : 24 = 10\text{ (m)}$.
}
\end{ex}

\begin{ex} %[5H5V26]
Một thửa ruộng hình thang có tổng độ dài hai đáy là $42\text{ m}$. Nếu mở rộng đáy lớn thêm $4\text{ m}$ thì diện tích thửa ruộng tăng thêm $30\text{ m}^2$. Tính diện tích thửa ruộng ban đầu theo đơn vị mét vuông.
\shortans{315}
\loigiai{
Chiều cao của thửa ruộng là:
$30 \times 2 : 4 = 15\text{ (m)}$.\\
Diện tích thửa ruộng ban đầu là:
$42 \times 15 : 2 = 315\text{ (m}^2\text{)}$.
}
\end{ex}

\begin{ex} %[5H5V26]
Một hình thang có đáy bé $12\text{ cm}$, đáy lớn gấp đôi đáy bé, chiều cao bằng trung bình cộng hai đáy. Diện tích của hình thang đó là bao nhiêu xăng-ti-mét vuông?
\shortans{324}
\loigiai{
Đáy lớn của hình thang là:
$12 \times 2 = 24\text{ (cm)}$.\\
Tổng độ dài hai đáy là:
$12 + 24 = 36\text{ (cm)}$.\\
Chiều cao hình thang là:
$36 : 2 = 18\text{ (cm)}$.\\
Diện tích hình thang là:
$36 \times 18 : 2 = 324\text{ (cm}^2\text{)}$.
}
\end{ex}

\begin{ex} %[5H5V26]
Một mảnh đất hình thang có đáy lớn $30\text{ m}$, đáy bé $20\text{ m}$. Người ta mở rộng đáy lớn thêm $6\text{ m}$ thì diện tích tăng thêm $45\text{ m}^2$. Chiều cao của mảnh đất đó bằng bao nhiêu mét?
\shortans{15}
\loigiai{
Chiều cao mảnh đất là:
$45 \times 2 : 6 = 15\text{ (m)}$.
}
\end{ex}

\begin{ex} %[5H5V26]
Cho hình thang vuông $ABCD$ có chiều cao $AD = 16\text{ cm}$, đáy bé $AB = 10\text{ cm}$ và đáy lớn $CD = 25\text{ cm}$. Diện tích của hình tam giác $BCD$ là bao nhiêu xăng-ti-mét vuông?
\shortans{200}
\loigiai{
Hình tam giác $BCD$ có độ dài cạnh đáy $CD = 25\text{ cm}$ và chiều cao hạ từ đỉnh $B$ xuống cạnh đáy $CD$ chính bằng chiều cao hình thang $AD = 16\text{ cm}$.\\
Diện tích hình tam giác $BCD$ là:
$25 \times 16 : 2 = 200\text{ (cm}^2\text{)}$.
}
\end{ex}
